\section{Limitations}
\label{sec:limitations}
The exact TimeSage-EV evaluated snapshot remains unavailable, so the BEA/NOAA/EIA replay is independent mechanistic evidence rather than an exact rerun. Its 35 endpoints are uneven across mechanisms: EIA is cutoff-focused; same-domain grounding transfer has only two held-out endpoints; cross-domain grounding is null; and held-out Kimi alignment has three independent endpoints. Broad cross-domain or model-invariant transfer is outside the claim.

Evidence layers have unequal status: DeepSeek is primary, Kimi is the support-repair replication with one provenance-invalid repeat excluded, the mini-model layer is exploratory, and EIA is post-review validation frozen before its outcomes. Repeats measure robustness, not task sample size. Pre-outcome family scorers remain primary; chronology, sensitivity analyses, intervals, and provenance appear in Appendix~A.

Two identification gaps remain open. The matched generic controls wrapper/complexity but is not guaranteed behaviorally neutral; a second neutral generic has not been executed. Nor have we executed a behaviorally equivalent retrieval-side temporal treatment needed for any container comparison. Both remain scientific-reopen conditions. Appendix~A closes narrower invocation, order, prompt-length, and ISO-date-serialization concerns; these diagnostics do not replace the missing baselines.

\section{Conclusion}
Reusable temporal skills become intervention-defined causal bottlenecks only in specific non-ceiling model--task regimes. A binding repair must beat both the original agent and matched generic assistance; this rule turns one apparent +28-point cutoff effect into no repair. Grounding, cutoff filtering, and alignment provide positive non-ceiling cases, while unchanged-code tests retain ceiling and cross-domain nulls. For self-evolving agents, recurring failure can nominate a persistent procedure, but promotion should require both controls; retrieval-side implementation remains a separate question.
